% Generated draft: requires mathematical and source review.
\documentclass[11pt]{article}
\usepackage[margin=25mm]{geometry}
\usepackage{fontspec}
\setmainfont{Latin Modern Roman}
\usepackage{amsmath,amssymb,mathtools,unicode-math}
\setmathfont{Latin Modern Math}
\usepackage{enumitem,needspace,xurl,hyperref}
\setlength{\parindent}{0pt}
\setlength{\parskip}{5pt}
\setlength{\emergencystretch}{3em}
\newcommand{\N}{\mathbb{N}}
\newcommand{\Z}{\mathbb{Z}}
\newcommand{\Q}{\mathbb{Q}}
\newcommand{\R}{\mathbb{R}}
\newcommand{\C}{\mathbb{C}}
\newcommand{\F}{\mathbb{F}}
\newcommand{\A}{\mathbb{A}}
\newcommand{\PP}{\mathbb{P}}
\newcommand{\E}{\mathbb{E}}
\newcommand{\eps}{\varepsilon}
\newcommand{\dd}{\,\mathrm{d}}
\begin{document}
\section*{UnsolvedMath ID: 20002059}
\textbf{Source problem code:} catalogue entry
\subsection{Definitions and mathematical statement}
This record is contextual prose between Problems 34 and 35, not a standalone open problem. The supplied verification restores its opening as ``Solutions to Problem 34 have a role to play'' in characterizing \(Q\)-curvature and GJMS operators.

Let \(M^n\) be a compact oriented even-dimensional manifold, not necessarily connected; assume no boundary for the stated integration identities. Write \(\Omega^k(M)\) for smooth forms and \(\mathcal C^k=\ker(\dd:\Omega^k\to\Omega^{k+1})\). The source describes generalized \(Q\)-operators \(T_g^k:\Omega^k\to\Omega^{n-k}\), for \(0\leq k\leq n/2-1\), with their intended weighted/Hodge identifications. A precise version of the supplied abstract descent lemma assumes
\[
T_{e^{2\omega}g}^k(c)=T_g^k(c)+L_k(\omega c)\qquad(c\in\mathcal C^k),
\quad L_k=\dd A_k\dd,
\]
where \(A_k:\Omega^{k+1}\to\Omega^{n-k-1}\), and \(L_k\) is conformally invariant and formally self-adjoint for the integrated form pairing. Multiplication by \(\omega\) occurs \emph{before} applying \(L_k\); it is not the scalar expression \((L_k\omega)c\).

The aid reports the resulting invariant space and map
\[
\mathcal H_g^k=\{c\in\mathcal C^k:\dd T_g^k(c)=0\},\qquad
c\longmapsto[T_g^k(c)]\in H_{\mathrm{dR}}^{n-k}(M).
\]
The space \(\mathcal H_g^k\) and this map are independent of representative, since \(\dd L_k=0\) and the change \(L_k(\omega c)\) is exact. For \(u\in\ker L_k\) and \(c\in\mathcal C^k\), the pairing
\[
\int_Mu\wedge T_g^k(c)
\]
is also conformally invariant by formal self-adjointness. The source additionally reports formal self-adjointness of the generalized \(Q\)-operators themselves and nontriviality of the induced map on some conformal manifolds; those properties are not inferred from factorization alone.

For \(k=0\), let \(q_g=Q_g\,\dd V_g\) and \(T_g^0(c)=c q_g\), with \(c\) locally constant. Its affine law is \(q_{e^{2\omega}g}=q_g+P_n\omega\), where the critical operator has the indicated exact-form factorization and kills locally constant functions. The lemma implies that \(\int_Mu c q_g\) is invariant for \(u\in\ker P_n\). Taking \(c\) to be each connected-component indicator shows that every component's total \(Q\)-curvature is separately invariant, not merely their sum.

The source attributes the geometric operator construction to Branson--Gover and writes \(L_k\) as a constant multiple of \(\dd M_{k+1}^g\dd\), with conventions to be supplied from their paper. At degree \(n/2-1\) it gives
\[
M^g=\dd\delta+2J-4P^{\sharp}:
\Omega^{n/2-1}\longrightarrow\Omega^{n/2-1}[-2],
\]
where \(\delta\) is the metric codifferential, \(J\) the Schouten trace, and \(P^{\sharp}\) its induced action on forms; conformal Hodge duality gives the complementary-degree version. It reports that these operators cannot generally be made conformally invariant merely by subtracting a divergence, invoking the leading Maxwell operator. This is retained as construction-specific background, not a consequence of the minimal lemma or a new unresolved question.
\subsection{Short English statement}
The record explains generalized \(Q\)-operators rather than posing an independent question. The supplied aid isolates an affine-law/factorization lemma yielding conformal-harmonic classes and nullspace pairings. In degree zero it gives separate conformal invariance of total \(Q\)-curvature on every connected component.
\subsection{Sources}
\begin{itemize}
\item Context between Problems 34 and 35, 2003 AIM workshop \emph{Conformal Structure in Geometry, Analysis, and Physics}: \url{https://aimath.org/WWN/confstruct/confstruct.pdf}.
\item T. Branson and R. Gover, supplied arXiv identifier \texttt{math.DG/0309085}, for the ambient/tractor generalized \(Q\)-operator construction. No full title or publication entry is included in the record.
\item Local registry \texttt{data/unsolvedmath/problems.json}, permanent ID 20002059, code \texttt{catalogue entry}; canonical \texttt{aim-geometry-notes.json}, \texttt{notes[396]}; updated 14 August 2026. The supplied verification identifies lost opening words and the contextual nature of the extract.
\item Attempt 1: context only, problem status at run not a problem, rigor label proved, novelty confidence low; current registry nevertheless says open/open triage. The machine-generated aid is not peer reviewed and requires independent verification before citation.
\end{itemize}
\end{document}
